\section{The Remaining Intellectual Virtues}
\label{sec:intellectual}

Aristotle identifies five ways by which the soul possesses truth without error in affirmation or denial: art, science, prudence, wisdom, and intellect (\emph{Nicomachean Ethics} VI.3). Prudence has already been treated because it is also cardinal and because it depends upon right appetite for its complete act. The other four perfect speculative reason or productive reason. Their measure is truth or the right work, not a quantitative mean of passion (ST I--II, q.~57).

\virtuedossier
  {Understanding of principles (\emph{nous}/\emph{intellectus})}
  {Principal intellectual virtue; Aristotle, \emph{Nicomachean Ethics} VI.6; ST I--II, q.~57, a.~2.}
  {Indemonstrable first principles that make demonstration and ordered inquiry possible.}
  {To apprehend a first principle as true and foundational rather than infer it from a prior demonstration.}
  {A \textbf{false denial} misses truth by deficiency at the level of judgment, but no dedicated passion-like vice habit flanks understanding (\oppositioncode{U}).}
  {A \textbf{false affirmation} misses truth by excess at the level of judgment (\oppositioncode{U}); multiplying axioms without warrant is faulty reasoning, not too much understanding.}
  {This acquired intellectual habit is distinct from the infused gift of understanding, from intelligence as measured aptitude, and from prudence's situational grasp.}

\virtuedossier
  {Science (\emph{episteme}/\emph{scientia})}
  {Principal intellectual virtue; Aristotle, \emph{Nicomachean Ethics} VI.3; ST I--II, q.~57, a.~2.}
  {Necessary or stable conclusions known through their causes by demonstration within a discipline.}
  {To demonstrate and assent to conclusions as following from true principles by a valid explanatory order.}
  {A \textbf{false denial} misses demonstrated truth by deficiency, while culpable ignorance can be a separate privation where knowledge is due; neither supplies a dedicated mean-flank vice habit (\oppositioncode{U}).}
  {A \textbf{false affirmation} misses truth by excess (\oppositioncode{U}); scientism is a philosophical overreach about method, not superabundant \emph{scientia}.}
  {The classical habit is not coextensive with modern natural science, accumulated information, credentials, or every probable conclusion. It can be morally misused because it perfects knowing rather than appetite.}

\virtuedossier
  {Wisdom (\emph{sophia}/\emph{sapientia})}
  {Highest principal speculative intellectual virtue by object; Aristotle, \emph{Nicomachean Ethics} VI.7; ST I--II, q.~57, a.~2; q.~66, a.~5.}
  {The highest causes, and lower truths judged and ordered in their light.}
  {To know the supreme causes available to the discipline and to judge the ordered whole through them.}
  {A \textbf{false denial} misses truth by deficiency at the level of judgment (\oppositioncode{U}); folly belongs to a different moral-theological analysis rather than a quantitative flank of this intellectual habit.}
  {A \textbf{false affirmation} misses truth by excess (\oppositioncode{U}); speculative pretension is counterfeit wisdom rather than wisdom beyond measure.}
  {Philosophical wisdom differs formally from the theological virtue of faith and from the Holy Spirit's gift of wisdom, though grace can order and elevate intellectual inquiry.}

\virtuedossier
  {Art (\emph{techne}/\emph{ars})}
  {Principal productive intellectual virtue; Aristotle, \emph{Nicomachean Ethics} VI.4; ST I--II, q.~57, aa.~3--4.}
  {A work capable of being made, considered according to the true plan and standards internal to its production.}
  {To conceive and execute the right form of a work through reliable productive reason.}
  {A \textbf{deficient false judgment} can fail to conform the productive plan to the work's rule (\oppositioncode{U}); incompetence can be nonculpable, while negligence belongs to a wider moral judgment.}
  {An \textbf{excessive false judgment} can posit what the work's true plan does not warrant (\oppositioncode{U}); virtuosity used for an evil end remains technically excellent while morally corrupted in its use.}
  {Art perfects the work as made; prudence perfects the maker's action as human. The fine arts are one subset of \emph{ars}, which also includes crafts and disciplined productive skills.}

\section{Qualified and Cognate States}
\label{sec:qualified}

The following realities have sound Aristotelian, Augustinian, or Thomistic provenance and are indispensable to a comprehensive study, but their own sources qualify their status. They are therefore not silently promoted into the strict virtue census. Their dossiers preserve the same analytical fields so that the reason for qualification is visible.

\virtuedossier
  {Natural or inchoate virtue}
  {Natural aptitude, seed, or likeness of virtue, not perfected moral virtue in the strict sense; Aristotle, \emph{Nicomachean Ethics} VI.13, 1144b1--17; ST I--II, q.~63, a.~1.}
  {Natural principles of reason and appetite, together with temperamental capacities that can incline toward acts resembling virtue.}
  {To supply an aptitude or beginning that right reason, prudence, choice, and habituation must form into acquired virtue; grace can heal and elevate the person beyond that natural proportion.}
  {No defect-side vice applies merely from weaker aptitude (\oppositioncode{NA}); temperament is not culpability.}
  {No excess-side vice applies merely from stronger aptitude (\oppositioncode{NA}); an unguided native tendency can become harmful because it lacks prudence, not because perfected virtue is excessive.}
  {Natural courage, sympathy, quick judgment, or restraint is neither moral destiny nor a settled habit. This dossier makes explicit the qualified state already distinguished in the philosophical grammar.}

\virtuedossier
  {Continence}
  {Praiseworthy state, broadly called a virtue but not perfect moral virtue in the strict Aristotelian sense; Aristotle, \emph{Nicomachean Ethics} VII.1--10; Augustine, \emph{On Continence}; ST II--II, qq.~155--156.}
  {Disordered strong desire that conflicts with a sound rational choice.}
  {To stand by right judgment and refuse the wrongful act despite appetite that has not yet been harmonized.}
  {\textbf{Incontinence} is the received named contrary in which judgment yields to desire (\oppositioncode{K}); because continence is qualified, it remains outside the strict table.}
  {No excess-side opposition applies (\oppositioncode{NA}); insensibility is a different disorder, and virtuous harmony is better than continued interior conflict.}
  {The continent person acts rightly but desires discordantly; the temperate person also has appetite ordered. Continence is morally valuable and often necessary in formation, yet it is not the completed harmony of virtue.}

\virtuedossier
  {Shame (\emph{aidos}/\emph{verecundia})}
  {Praiseworthy passion and disposition, not a virtue strictly speaking; Aristotle, \emph{Nicomachean Ethics} II.7 and IV.9; ST II--II, q.~144.}
  {Disgrace apprehended as an evil attached to a base action.}
  {To recoil affectively from disgrace in a manner that can deter wrongdoing or assist repentance.}
  {\textbf{Shamelessness} has deficient fear of deserved disgrace (\oppositioncode{M}, qualified because the mean is a passion).}
  {\textbf{Bashfulness} fears disgrace excessively or where no base act warrants it (\oppositioncode{M}, qualified).}
  {The mature virtuous person chiefly avoids base action because it is base, not merely from fear of exposure. Toxic humiliation and involuntary embarrassment are not moral virtue.}

\virtuedossier
  {Complete friendship}
  {Aristotle calls friendship a virtue or something involving virtue; the complete kind is constituted by reciprocal love of the good, \emph{Nicomachean Ethics} VIII.1, 3--4 and IX.4, 9; ST II--II, q.~23, a.~1; q.~114, a.~1. It is treated as a virtuous relation rather than one additional habit.}
  {The friend's good and a shared life in which each loves the other as another self for the other's character and good.}
  {To will and communicate good reciprocally, live together in fitting activity, and sustain truthful concord.}
  {No single defect-side vice exhausts failed friendship (\oppositioncode{NA}); isolation, exploitation, betrayal, and indifference differ formally.}
  {No excess-side opposition applies (\oppositioncode{NA}); possessiveness is a corruption of love and justice, not friendship perfected too far.}
  {Friendships of utility and pleasure are real but incomplete and fragile. Complete friendship differs from affability toward ordinary associates and from supernatural charity's friendship with God.}

\virtuedossier
  {Righteous indignation (\emph{nemesis})}
  {Praiseworthy mean passion named by Aristotle, not a settled coordinate virtue in the developed census; \emph{Nicomachean Ethics} II.7, 1108a30--b6; ST II--II, q.~36, aa.~1--2.}
  {The apparent distribution of good fortune in relation to desert.}
  {To feel proportionate pain at undeserved prosperity because one loves justice, without willing evil to the prosperous person.}
  {\textbf{Spite} is deficient in the relevant pain and can rejoice at another's misfortune (\oppositioncode{M}, qualified).}
  {\textbf{Envy} is pained at another's good as such or beyond desert's relevance (\oppositioncode{M}, qualified).}
  {The passion does not authorize private punishment or certainty about hidden desert. Christian joy in another's good and grief over injustice require charity, prudence, and epistemic humility.}

\virtuedossier
  {Heroic or divine virtue}
  {Exceptional degree or mode set over against brutishness; Aristotle, \emph{Nicomachean Ethics} VII.1, 1145a15--33; ST I--II, q.~68, a.~1 ad 1; II--II, q.~159, a.~2 ad 1. It is not a further species parallel to courage or justice.}
  {The surpassing stability and excellence with which the whole rational life may embody moral good.}
  {To perform the acts of the virtues with a rare constancy and elevation that appears more than ordinary human excellence.}
  {\textbf{Bestiality} is Aristotle's named qualitative contrary (\oppositioncode{K}), not a deficient quantity of ordinary virtue; ordinary genuine virtue is not vice.}
  {No excess-side opposition applies (\oppositioncode{NA}); the heroic names excellence of mode, not a quantity that can overshoot reason.}
  {The Aristotelian category does not by itself establish canonized heroic virtue, infused grace, sanctity, miracles, or a juridical cause of saints. Those later theological and canonical judgments require their own sources.}

\virtuedossier
  {Nobility-and-goodness (\emph{kalokagathia})}
  {Integrative Aristotelian capstone especially explicit in \emph{Eudemian Ethics} VIII.3; retained as a qualified whole rather than a coordinate virtue.}
  {The noble exercise and ordering of the virtues and external goods under the complete human good.}
  {To choose noble acts for their own sake within an integrated excellent life.}
  {No separate defect-side vice applies (\oppositioncode{NA}); vice and fragmentation name many failures rather than one paired extreme.}
  {No separate excess-side vice applies (\oppositioncode{NA}); it is a title for integrated excellence.}
  {The term is not a license to import aristocratic birth, appearance, social class, or modern prestige into moral worth. It does not replace the distinct virtues whose integration it names.}

\subsection{Why qualification matters}

A passion can be fitting without being a habit; a relation can involve virtue without being another power's perfection; a state can resist vice without giving appetite its final harmony; and an integrative title can name the whole without adding one part. Counting each as a coordinate virtue would make the catalogue larger but the philosophy poorer. Omitting them would conceal exactly the distinctions needed for formation.
